\documentclass[11pt,a4paper]{article}
\usepackage[utf8]{inputenc}
\usepackage[T1]{fontenc}
\usepackage{lmodern}
\usepackage{textcomp}
\usepackage[spanish,es-noquoting]{babel}
\usepackage{amsmath,amssymb}
\usepackage{booktabs}
\usepackage{longtable}
\usepackage{geometry}
\geometry{margin=2.4cm}
\usepackage{enumitem}
\usepackage{xcolor}
\usepackage{hyperref}
\hypersetup{colorlinks=true,linkcolor=blue!55!black,urlcolor=blue!55!black,citecolor=blue!55!black}
\usepackage{fancyhdr}
\pagestyle{fancy}\fancyhf{}\rhead{SCRES-IA — la carrera completa}\lhead{\thepage}\cfoot{}
\setlength{\headheight}{14pt}
\newcommand{\ret}{\text{ReT}}
\newcommand{\code}[1]{\texttt{#1}}

\title{\textbf{SCRES-IA: la carrera completa}\\[4pt]
\large Resiliencia, aprendizaje por refuerzo y control estructurado\\
en una cadena de suministro militar de alimentos\\[8pt]
\normalsize Documento acompañante de \emph{Submission A} (Program Q) y del Paper 2}
\author{Programa de investigación SCRES-IA}
\date{22 de julio de 2026}

\begin{document}
\maketitle
\thispagestyle{fancy}

\begin{abstract}
\noindent Este documento explica, de principio a fin y sin dar nada por supuesto, qué es el programa
SCRES-IA, qué hemos logrado, qué significa cada término técnico que usamos y en qué punto exacto de la
carrera estamos. Está pensado como \emph{compañero pedagógico} de dos entregables: (i) \emph{Submission
A} / Paper 2, el manuscrito sobre \emph{Program Q} que ya tenemos; y (ii) el frente activo de
aprendizaje acumulado entre campañas (\emph{Q-R1}). No sustituye a los manuscritos: los explica. El
lector no necesita saber de antemano qué es un ``lazo abierto'' ni un ``MPC''; todo se define aquí. La
tesis central, ya sostenida por evidencia sellada, es doble y honesta: \textbf{el aprendizaje por
refuerzo adquiere control adaptativo real —le gana a las 65\,536 políticas fijas posibles— pero no
supera al control estructurado clásico}; y \textbf{la retención de conocimiento entre campañas mejora
causalmente la resiliencia inicial}, aunque todavía no de forma conjuntamente segura.
\end{abstract}

\tableofcontents
\bigskip
\hrule
\bigskip

%==============================================================
\section{Para qué sirve este documento}
El programa SCRES-IA (\emph{Supply-Chain RESilience with Inteligencia Artificial}) ha producido muchos
experimentos, veredictos y correcciones. Este texto los ordena en una sola narrativa. Responde tres
preguntas que se hicieron explícitamente:
\begin{enumerate}[nosep]
  \item \textbf{¿Qué logramos?} (Secciones \ref{sec:logros} y \ref{sec:qr1}.)
  \item \textbf{¿Qué significa cada término?} (Sección \ref{sec:glosario}, un glosario exhaustivo:
        lazo abierto/cerrado, MPC, RL, PPO, ReT, \emph{belief}, \emph{cold-start}, placebos,
        \emph{guardrails}, etc.)
  \item \textbf{¿Cómo fue la carrera y dónde estamos?} (Secciones \ref{sec:metodo}, \ref{sec:carrera}
        y \ref{sec:estado}.)
\end{enumerate}

%==============================================================
\section{El problema y el mundo simulado}\label{sec:mundo}

\subsection{La tesis de Garrido-Ríos (2017)}
El punto de partida es una tesis doctoral que modela una \textbf{cadena de suministro militar de
alimentos}: el flujo de raciones desde su producción hasta las tropas desplegadas, a través de
\textbf{13 operaciones} encadenadas (producción, procesamiento, almacenamiento, transporte por
distintos escalones, etc.). El interés no es el costo ni la ganancia, sino la \textbf{resiliencia}: la
capacidad de la cadena de seguir sirviendo pedidos a tiempo cuando ocurren disrupciones (averías,
retrasos de proveedores, picos de demanda, pérdidas de capacidad). La tesis define una métrica formal
de resiliencia, \textbf{ReT}, y la calcula en hojas de cálculo Excel a partir de datos de la operación.

\subsection{La petición de Garrido et al.\ (2024)}
Años después, un artículo de seguimiento plantea una agenda: los modelos de simulación tradicionales
\emph{no aprenden de la experiencia}. Cada corrida empieza de cero; el sistema no acumula sabiduría de
una campaña a la siguiente. Los autores lo llaman, medio en broma, el ``efecto Alzheimer'' de la
simulación, y proponen usar \textbf{inteligencia artificial} para dos cosas:
\begin{enumerate}[nosep]
  \item \textbf{Convertir la simulación de lazo abierto en lazo cerrado} (definidos abajo): que un
        controlador \emph{reaccione} al estado del sistema en vez de seguir un plan fijo.
  \item \textbf{Acumular aprendizaje entre iteraciones}: que la experiencia pasada mejore las
        decisiones futuras. Esta es la ``estrella polar'' del programa.
\end{enumerate}

\subsection{Nuestra reconstrucción y su fidelidad}
Reconstruimos la cadena de 13 operaciones como una \textbf{simulación de eventos discretos} (DES; ver
glosario) en Python/SimPy, y reprodujimos la métrica ReT de la tesis \textbf{exactamente}: sobre
$47\,546$ celdas de fórmula de los libros Excel reales, \textbf{0 discrepancias}. Esta fidelidad
numérica es el cimiento: cualquier cosa que midamos después descansa sobre un modelo verificado contra
la fuente.

\subsection{La extensión de dos productos: el problema de decisión}\label{sec:extension}
Sobre esa base construimos una \textbf{extensión de investigación}: dos productos no fungibles
(llamémoslos $P_C$ y $P_H$) que comparten la misma línea de procesamiento (operaciones Op5--Op7). Cada
semana, durante 8 semanas, hay que decidir \textbf{cómo repartir la producción entre los dos
productos}. Cada decisión semanal elige uno de 4 niveles de mezcla; con 8 decisiones, el espacio
completo de \emph{planes fijos} tiene
\[
4^{8} = 65\,536 \text{ calendarios posibles.}
\]
Este número es importante: es pequeño lo suficiente para \textbf{enumerarlo entero}. Podemos calcular
el ReT de \emph{cada uno} de los 65\,536 planes fijos y así conocer, exactamente, cuál es el mejor plan
que no reacciona al estado. Eso nos da una \emph{clave de respuestas} (``answer key'') contra la cual
medir si reaccionar al estado aporta algo.

%==============================================================
\section{Glosario exhaustivo}\label{sec:glosario}
Aquí se define cada término. Los nombres técnicos suelen quedarse en inglés (así aparecen en la
literatura y en el código); se explica su significado.

\begin{description}[style=unboxed,leftmargin=0pt]
\item[Simulación de eventos discretos (DES, \emph{Discrete-Event Simulation}).]
Un modelo de computadora donde el tiempo avanza ``saltando'' de un evento al siguiente (llega un
pedido, termina un lote, sale un camión), en vez de continuamente. Es la forma estándar de simular
cadenas de suministro. La nuestra usa la librería SimPy.

\item[Lazo abierto (\emph{open loop}).]
Un plan de acciones fijado \emph{de antemano} que se ejecuta pase lo que pase, sin mirar cómo va el
sistema. Analogía: un temporizador de riego que abre el agua a las 6\,a.m.\ todos los días, llueva o
truene. En nuestro problema, un ``calendario'' de 8 decisiones semanales fijadas antes de empezar es
una política de lazo abierto. Hay $65\,536$ de ellas.

\item[Lazo cerrado (\emph{closed loop}).]
Un controlador que \emph{observa el estado} del sistema y \emph{reacciona}: ajusta la próxima decisión
según lo que está pasando (inventario, pedidos atrasados, demanda observada). Analogía: un termostato
que mide la temperatura y enciende o apaga la calefacción. Convertir el DES de lazo abierto en lazo
cerrado es la primera petición de Garrido 2024.

\item[Resiliencia ReT (\emph{Resilience over Time}).]
La métrica central, tomada de la tesis. Mide, pedido por pedido, qué tan bien la cadena entrega a
tiempo. Sin disrupción activa, para el pedido $j$:
\[
\ret_j = 1 - \frac{B_{t}+U_{t}}{j},
\]
donde $B_t$ son pedidos en \emph{backorder} (atrasados pero aún en cola) y $U_t$ pedidos no atendidos.
Con disrupción activa hay ramas adicionales (autotomía, recuperación). Un ReT de 1 es servicio
perfecto; valores menores indican degradación. \textbf{Nunca modificamos esta fórmula}: es fiel a la
tesis (0/47\,546).

\item[MPC (\emph{Model Predictive Control}, control predictivo por modelo).]
Un controlador \emph{estructurado clásico} (no una red neuronal). En cada decisión: (1) usa un modelo
del sistema para simular hacia adelante un \emph{horizonte} de varias semanas; (2) evalúa muchas
secuencias de acciones candidatas; (3) elige la primera acción de la mejor secuencia; (4) avanza una
semana y repite. Es ``receding-horizon'' (horizonte deslizante). Es fuerte, interpretable y respeta
restricciones de seguridad. En este trabajo, el rival a vencer.

\item[\emph{belief} / filtro bayesiano.]
La ``creencia'' del controlador sobre variables ocultas del mundo (por ejemplo, cuál producto va a
dominar la demanda). Un \emph{filtro bayesiano} actualiza esa creencia de forma óptima a medida que
llegan observaciones. Un \textbf{belief-MPC} es un MPC que planifica sobre esa creencia. Punto técnico
clave: para un régimen de dos estados con dinámica conocida, la creencia bayesiana es un
\emph{estadístico suficiente} —contiene toda la información útil de la historia—, así que \emph{ninguna
red puede formar una mejor creencia} a partir de los mismos datos.

\item[Aprendizaje por refuerzo (RL, \emph{Reinforcement Learning}).]
Una familia de métodos donde un agente \emph{aprende} una política (qué hacer en cada estado)
maximizando una recompensa, por ensayo y error sobre muchos episodios simulados. A diferencia del MPC,
no planifica en línea: aprende una política y luego la ejecuta con un solo ``forward pass''.

\item[PPO / RecurrentPPO.]
\emph{Proximal Policy Optimization}: el algoritmo de RL que usamos. \emph{RecurrentPPO} le añade
memoria (una red LSTM) para que la política pueda recordar lo observado dentro de un episodio. Es
nuestro aprendiz principal.

\item[\emph{cold-start} (arranque en frío).]
Las primeras decisiones de una campaña, \emph{antes} de haber observado casi nada de esa campaña. Es
donde el conocimiento retenido de campañas anteriores debería ayudar más: uno decide mejor al inicio
si recuerda cómo fueron las campañas pasadas.

\item[\emph{retained} vs \emph{reset} (retenido vs reiniciado).]
El tratamiento central de Q-R1. Tras terminar una campaña se hace un \textbf{reset físico completo}
(inventario, colas, pedidos en tránsito, todo a cero). La única diferencia entre brazos es si el
\emph{conocimiento} (la creencia/posterior) se \textbf{retiene} de la campaña anterior o se
\textbf{reinicia}. Si el brazo retenido decide mejor en frío que el reiniciado, ese es el valor causal
del conocimiento acumulado.

\item[Cohorte temprana y \emph{full-cohort} vs \emph{visible}.]
Medimos la ReT sobre la \emph{cohorte} de pedidos generados en las 2 primeras semanas. El endpoint
\textbf{visible} promedia solo los pedidos completados (el defecto que las auditorías señalaron: puede
``ganar'' dejando pedidos difíciles sin resolver). El endpoint \textbf{full-cohort} mantiene los 12
pedidos de la cohorte en el denominador, puntuando 0 a los no resueltos o perdidos. Usamos el
full-cohort como endpoint honesto.

\item[\emph{worst-product fill} (peor producto).]
El nivel de servicio del \emph{producto peor atendido}. Es un \emph{guardrail} de equidad: impide
subir el ReT medio favoreciendo a un producto y matando de hambre al otro.

\item[\emph{unresolved} / \emph{lost} / \emph{backorder}.]
Pedidos \emph{no resueltos} (aún sin atender al corte), \emph{perdidos} (expulsados de la cola) y
\emph{atrasados} (en cola). Guardrails de servicio: una mejora de ReT no cuenta si aumenta estos.

\item[Dosis-respuesta (\emph{dose-response}).]
Prueba de mecanismo. Variamos la \emph{persistencia} $\kappa$ entre campañas (qué tanto se parece una
campaña a la anterior). Si el efecto retenido es real, debe \emph{crecer con la persistencia}:
$\kappa{=}0.90 > \kappa{=}0.75 > \text{iid}$ (iid $=$ campañas independientes $=$ nada que recordar).

\item[Placebos (\emph{placebo controls}).]
Brazos de control negativo que \emph{no deberían} producir mejora si el mecanismo es genuino:
\emph{shuffled} (historia barajada), \emph{wrong} (historia deliberadamente equivocada, un
``anti-oráculo''), \emph{delayed} (historia retrasada una campaña). Si un placebo mejora, hay una fuga
en el mecanismo.

\item[\emph{guardrail}.]
Una barrera de validez que se debe pasar aparte de la magnitud del efecto: recursos exactamente
iguales, sin aumento de perdidos/no-resueltos, equidad de producto, misma información para ambos
brazos, reset físico idéntico. Sin ellas, ``más ReT'' se puede comprar haciendo trampa.

\item[SESOI y LCB (\emph{Smallest Effect Size Of Interest}; \emph{Lower Confidence Bound}).]
SESOI es el tamaño de efecto mínimo que nos importa científicamente ($+0.01$ de ReT). LCB es el
extremo inferior de un intervalo de confianza al 95\%: exigimos LCB${}>0$ (o ${}\ge$ SESOI) para que el
efecto no sea ruido. Los intervalos son \emph{clustered bootstrap} por historia (respetan que los
pedidos de una misma historia están correlacionados).

\item[CRN (\emph{Common Random Numbers}).]
Técnica para comparar dos políticas de forma justa: se les aplica \emph{exactamente el mismo azar}
(los mismos shocks de demanda), de modo que las diferencias se deban a la política, no a la suerte.

\item[TOST / equivalencia.]
\emph{Two One-Sided Tests}: prueba estadística para afirmar que dos cosas son \emph{prácticamente
iguales} (no solo que ``no encontramos diferencia''). Así probamos rigurosamente RL${\approx}$MPC.

\item[Las cuatro cantidades de la escalera: $H_{PI}$, $H_{obs}$, $H_{OL}$, $\Delta_N$.]
El corazón metodológico (Sección \ref{sec:metodo}):
\begin{itemize}[nosep]
  \item $H_{PI}$ — \emph{oportunidad física} clarividente: cuánto ReT ganaría un oráculo que conoce el
        futuro. Es el techo.
  \item $H_{obs}$ — \emph{conversión observable}: cuánto de esa oportunidad convierte un controlador
        \emph{clásico} desplegable (sin ver el futuro).
  \item $H_{OL}$ — \emph{adaptación aprendida}: cuánto le gana el RL a la \emph{mejor política fija}
        (las 65\,536). Mide si reaccionar al estado sirve.
  \item $\Delta_N$ — \emph{prima neural}: cuánto le gana el RL al \emph{mejor MPC estructurado}. Es la
        pregunta ``¿vale la pena la red?''.
\end{itemize}

\item[Seeds selladas vs quemadas (\emph{sealed} vs \emph{burned}).]
Una \emph{seed} fija el azar de una corrida. Las \emph{quemadas} son de desarrollo (se pueden mirar
muchas veces). Las \emph{selladas} se abren \emph{una sola vez} para la confirmación prospectiva: se
congela el contrato y los criterios \emph{antes} de abrirlas, para que nadie ``pinte la diana después
del disparo''. Disciplina anti-autoengaño.

\item[Custodia y A11.]
Cada resultado guarda hashes SHA-256, commit, entorno y filas crudas. ``A11'' es nuestra regla:
\emph{verificar cada afirmación contra el artefacto real} antes de creerla —incluso las de nuestro
propio agente colaborador.
\end{description}

%==============================================================
\section{El método: la escalera de falsación de 4 niveles}\label{sec:metodo}
La contribución metodológica de \emph{Submission A} no es un algoritmo de RL nuevo, sino una
\textbf{forma de decidir, con rigor, cuándo un controlador aprendido aporta valor sobre el control
estructurado}. Se organiza como una escalera que \emph{localiza dónde vive (o muere) el valor}:
\[
\underbrace{H_{PI}}_{\text{¿hay oportunidad física?}}
\;\to\;
\underbrace{H_{obs}}_{\text{¿la convierte el control clásico?}}
\;\to\;
\underbrace{H_{OL}}_{\text{¿el RL aprende a adaptarse?}}
\;\to\;
\underbrace{\Delta_N}_{\text{¿la red supera a la estructura?}}
\]
Cada peldaño es \emph{falsable por separado}. La clave es la \textbf{clave de respuestas exacta}: como
enumeramos las 65\,536 políticas fijas, sabemos exactamente cuál es la mejor política que no reacciona.
Así, cuando el RL le gana a \emph{todas}, es adaptación genuina y no ``redescubrir un buen plan fijo''
—una distinción que los trabajos de RL-en-DES con baselines débiles no pueden hacer.

%==============================================================
\section{Lo que logramos: Submission A / Paper 2 (Program Q)}\label{sec:logros}

\subsection{El entorno y el diseño}
Demanda con cambio de régimen (dos productos), 8 decisiones semanales, 4 acciones por decisión,
$65\,536$ calendarios de lazo abierto enumerados, \textbf{sin riesgos activos} (entorno estacionario),
métrica ReT canónica. Confirmación con $256$ ``tapes'' (realizaciones de azar) por celda,
$21\,696$ replays directos del DES completo, \textbf{0 fallos de auditoría}, 3 celdas de régimen.

\subsection{Resultado 1 — el RL adquiere adaptación real}
Una política recurrente (RecurrentPPO), entrenada desde cero, \textbf{le gana a las 65\,536 políticas
fijas} en las tres celdas, con márgenes grandes de ReT canónica ($H_{OL}$):
\begin{center}
\begin{tabular}{lccc}
\toprule
Celda de régimen & $H_{OL}$ (RL $-$ mejor fija) & Semillas positivas \\
\midrule
rho75\_share90 & $+0.0795$ & 10/10 \\
rho90\_share75 & $+0.0726$ & 10/10 \\
rho90\_share90 & $+0.1172$ & 10/10 \\
\bottomrule
\end{tabular}
\end{center}
Esto \textbf{entrega la primera petición de Garrido 2024}: el DES de lazo abierto convertido en lazo
cerrado por una IA que \emph{aprende a reaccionar}, medido con rigor. La curva de aprendizaje sube
monótonamente y converge \emph{desde abajo} al nivel del control estructurado, sin deriva.

\subsection{Resultado 2 — pero no hay prima neural: RL${\approx}$MPC}
Cuando comparamos el RL con el \emph{mejor control estructurado} (belief-MPC), la diferencia
$\Delta_N$ es \textbf{estadísticamente equivalente a cero} (TOST) en las tres celdas:
\begin{center}
\begin{tabular}{lcc}
\toprule
Celda & $\Delta_N$ (RL $-$ MPC) & IC 95\% \\
\midrule
rho75\_share90 & $-0.0016$ & $[-0.0063,\,+0.0031]$ \\
rho90\_share75 & $-0.0007$ & $[-0.0055,\,+0.0041]$ \\
rho90\_share90 & $-0.0004$ & $[-0.0027,\,+0.0019]$ \\
\bottomrule
\end{tabular}
\end{center}
Los intervalos \emph{cruzan y encierran} el cero: no es ``no encontramos diferencia'', es
\emph{equivalencia demostrada}. La igualdad sobrevivió a \textbf{cinco pruebas adversariales
independientes} (replicación prospectiva, una arquitectura distinta tipo transformer, arranque
tibio desde tres inicializaciones, un factorial de geometría de recompensa, y un aprendiz off-policy
QR-DQN). Ningún ``defecto arreglable'' del RL explica el empate.

\subsection{Resultado 3 — el mecanismo: la decisión es \emph{belief-insensitive}}
¿Por qué empatan? Porque en este entorno estacionario \textbf{la decisión óptima casi no depende de la
creencia}: una mejor inferencia cambia la acción óptima en ${<}5\%$ de los estados, y la cima del
paisaje de políticas fijas es \emph{plana} (las 38 mejores dentro de $0.0025$). El control
estructurado ya se sienta en la frontera observable; no queda residual que una red pueda capturar.

\subsection{Resultado 4 — la amortización tampoco favorece a la red}
Se creía que aunque empataran en calidad, el RL podría ganar en \emph{costo}: un ``forward pass'' en
vez de planificar en línea. El benchmark de latencia (hardware-específico, 5\,000 repeticiones) lo
\textbf{refuta}:
\begin{center}
\begin{tabular}{lcc}
\toprule
Controlador & Mediana (ms) & p95 (ms) \\
\midrule
RecurrentPPO & $0.573$ & $0.702$ \\
Familia estructurada & $0.082$ & $0.088$ \\
\bottomrule
\end{tabular}
\end{center}
El control estructurado es \textbf{${\sim}7\times$ más rápido}. Así que la familia estructurada iguala
al RL en resiliencia (equivalencia práctica) \emph{y} lo supera en cómputo. Honestamente, entonces,
Submission A \textbf{no reclama ninguna ventaja de amortización neural}.

\subsection{El reencuadre: un marco de decisión, no un ``RL ganó''}
El veredicto congelado de Program Q es \code{STOP\_Q\_NO\_REPLICATED\_LEARNED\_ADAPTATION}: el RL
aprende la adaptación, pero \emph{no} hay prima sobre la estructura. Esto se presenta no como una
decepción sino como el \textbf{poder discriminante del marco}: la escalera separa ``hubo aprendizaje''
(Nivel 3, $+0.07$ a $+0.12$ sobre toda la frontera) de ``el aprendizaje le ganó a la estructura''
(Nivel 4, ${\approx}0$) —una distinción que los papers de baseline débil no pueden trazar. El título de
trabajo lo captura:
\begin{quote}
\emph{``When Feedback Beats an Exhaustive Static Frontier but Not Structured Control: Exact Benchmarking
of Recurrent RL in a Supply-Chain DES.''}
\end{quote}
Es un resultado más honesto y, paradójicamente, más fuerte: un \textbf{protocolo para decidir cuándo un
controlador aprendido merece desplegarse}, aplicado a un DES industrial validado. Destino: \emph{Computers
\& Industrial Engineering} (Q1).

%==============================================================
\section{El frente activo: Q-R1, el aprendizaje acumulado entre campañas}\label{sec:qr1}

\subsection{La estrella polar de Garrido}
La \emph{segunda} petición de Garrido 2024 —la que de verdad importa— es el \textbf{aprendizaje
acumulado}: que la experiencia de campañas pasadas mejore las decisiones de campañas futuras, después
de un reset físico completo. Ese es el ``efecto Alzheimer'' curado. Q-R1 es nuestro ataque a esta
pregunta.

\subsection{El diseño}
Se generan \emph{historias} de campañas correlacionadas (persistencia $\kappa$). Entre campañas, reset
físico total; solo el \emph{conocimiento} (posterior bayesiano) se retiene o se reinicia. Se mide la
ReT \emph{full-cohort} de la cohorte de arranque en frío. Brazos: retenido, reiniciado, y placebos
(shuffled, wrong, delayed). Comparador: la familia MPC estructurada más fuerte que congelamos y
verificamos (H4, integración estratificada, con \emph{gate} de convergencia p16/p64).

\subsection{El veredicto verificado (confirmación prospectiva, seeds selladas)}
La confirmación corrió ${\sim}9.3$\,h en el VPS sobre seeds selladas frescas (roots
$7\,572\,001$--$7\,572\,032$, commit congelado, integridad verificada A11). Verdicto:
\code{STOP\_REPAIRED\_Q\_R1\_NO\_RETAINED\_INFORMATION\_PASS}. Pero \textbf{la estructura del STOP es la
noticia}:
\begin{center}
\begin{tabular}{lc}
\toprule
Cantidad (κ$=0.90$, full-cohort) & Valor \\
\midrule
ReT retenido $-$ reiniciado & $\mathbf{+0.0656}$ (LCB95 $+0.0524$) \\
Dosis-respuesta & $+0.0656\;(\kappa.90) > +0.0341\;(\kappa.75) > 0\;(\text{iid})$ \\
Anti-oráculo (\emph{wrong}) & $-0.1016$ \\
\midrule
\emph{Falla} — placebo \emph{delayed} & $+0.0059\;(>0)$ \\
\emph{Falla} — worst-product (LCB) & $-0.0365\;(<-0.02)$; no-resueltos $+0.094$ \\
\bottomrule
\end{tabular}
\end{center}

\paragraph{Lectura honesta.} El efecto retenido es \textbf{real, grande y con dosis-respuesta
perfecta}, y el anti-oráculo daña fuertemente —exactamente la firma de un mecanismo causal genuino.
\textbf{El aprendizaje acumulado de Garrido queda confirmado prospectivamente} en seeds selladas, y
contra el comparador estructurado \emph{más fuerte} (el efecto es aún mayor que el $+0.0226$ de los
diagnósticos de desarrollo). \emph{Pero} no pasa dos barreras: (i) el placebo \emph{delayed} es
levemente positivo —una fuga menor de temporización—; y, sobre todo, (ii) el brazo retenido sube el
ReT medio \emph{en parte favoreciendo al producto mayoritario}, degradando el servicio del minoritario
por debajo del margen de equidad. Es la \textbf{misma forma que hallamos en Program O}: la conversión
\emph{media} es real; la conversión \emph{conjuntamente segura} no está establecida.

\subsection{Qué abre esto}
El STOP es correcto y \textbf{no autoriza entrenar ninguna red}. Pero define con precisión la próxima
pregunta: el ``residual convertible'' aquí \emph{no es de magnitud} (la magnitud sobra: $+0.066$), sino
de \textbf{seguridad conjunta}. La pregunta defendible es: \emph{¿puede un controlador retenido
consciente de restricciones —o un híbrido learning-augmented MPC— mantener la ganancia media de
$+0.066$ SIN el costo de peor-producto/no-resueltos?} Si un híbrido con \emph{gate} de confianza captura
esa ganancia respetando la equidad, ahí está el paper fuerte.

%==============================================================
\section{La carrera honesta: el ``cuándo NO entrenar''}\label{sec:carrera}
Una parte esencial del programa es lo que \emph{no} funcionó, documentado con la misma disciplina que
lo que sí. A lo largo de la carrera abrimos y cerramos muchos frentes (variantes de riesgo, prevención,
reordenamiento, perecederos, timing endógeno, atlas de estrés de guerra), y casi todos terminaron en
\textbf{STOP} bien fundamentados: había ``oportunidad clarividente'' ($H_{PI}>0$) pero \emph{no} se
convertía en ventaja observable, o el controlador estructurado ya capturaba todo el valor. El patrón
recurrente:
\begin{quote}
\emph{El aprendizaje por refuerzo no está justificado cuando el control estructurado ya se sienta en la
frontera observable. Saber cuándo NO entrenar es, en sí mismo, un resultado.}
\end{quote}
Esta honestidad tiene tres pilares metodológicos, todos vinculantes:
\begin{enumerate}[nosep]
  \item \textbf{Seeds selladas de un solo tiro} con contrato y criterios congelados \emph{antes} de
        abrirlas (no se puede pescar el resultado favorable).
  \item \textbf{Verificación adversarial (A11)}: cada afirmación se comprueba contra el artefacto
        real. Dos agentes se auditan mutuamente; varias sobre-afirmaciones (mías y del colaborador)
        fueron atrapadas y corregidas —incluida esta misma semana.
  \item \textbf{Métrica intacta}: la ReT canónica nunca se cambia (0/47\,546); las mejoras deben
        pasar guardrails de recursos, servicio y equidad, no comprarse con trampa.
\end{enumerate}

%==============================================================
\section{Estado actual y ruta}\label{sec:estado}
Trabajamos en dos carriles paralelos.

\paragraph{Carril A — publicar Submission A / Paper 2 (Program Q).} Listo en lo esencial: el marco de 4
niveles, la frontera exacta, la equivalencia RL${\approx}$MPC robusta a 5 pruebas, el mecanismo
belief-insensitive y el benchmark de latencia (que retira honestamente el argumento de amortización).
Destino \emph{Computers \& Industrial Engineering}. Es la publicación asegurada antes de la graduación.

\paragraph{Carril B — Q-R1 y el intento de ganarle limpio al MPC.} La secuencia disciplinada,
acordada tras varias rondas de auto-crítica adversarial, es:
\begin{enumerate}[nosep]
  \item Adjudicar la confirmación canónica de Q-R1 (\textbf{hecho}: STOP con efecto real pero no
        conjuntamente seguro).
  \item Congelar el mejor comparador estructurado \emph{probado} (nunca llamarlo ``MPC óptimo'').
  \item Medir el residual con full-cohort, CRN independiente de la política, restricciones completas y
        auditoría de potencia con SESOI $+0.01$ —\emph{sin umbrales arbitrarios}.
  \item Solo si queda residual: probar un \emph{learning-augmented MPC} (valor terminal aprendido y/o
        corrección de creencia), con RecurrentPPO puro como \emph{ablación}, no como candidato.
  \item Si no hay residual estacionario: sensibilidad \emph{learner-blind} de los riesgos existentes;
        y solo si son inertes, una extensión de riesgo acoplado a producto \textbf{validada
        físicamente por Garrido}, con física y grilla congeladas antes de cualquier RL.
\end{enumerate}
La apuesta más prometedora \emph{no} es más ajuste de PPO estacionario, sino \textbf{un encoder de
contexto causal retenido $+$ un MPC factible, bajo riesgos persistentes específicos por producto
físicamente validados} —donde una red puede aprender una representación más rica que el filtro discreto
actual, siempre que el MPC reciba exactamente la misma historia.

\paragraph{Qué significa ``ganarle al MPC'', con honestidad.} No existe un teorema de imposibilidad ni
un MPC óptimo certificado. ``Ganarle'' significa ganarle a \emph{la familia MPC retenida más fuerte que
hayamos congelado y probado} bajo el mismo contrato de información, servicio y recursos —y el candidato
final sería, él mismo, un \emph{learning-augmented MPC}. Sabemos que existe una \emph{incertidumbre
física históricamente aprendible} (el efecto retenido de $+0.066$); no sabemos aún si un controlador
justo deja un residual convertible \emph{y seguro}. Esa medición, no más PPO, es el próximo paso.

%==============================================================
\section{Limitaciones honestas}\label{sec:limitaciones}
\begin{itemize}[nosep]
  \item \textbf{Entorno estacionario y extensión de investigación.} La equivalencia RL${\approx}$MPC se
        demostró sin riesgos activos y con una extensión de dos productos \emph{definida por nosotros},
        no estimada de la cadena militar real. Bajo riesgos físicamente relevantes la respuesta podría
        cambiar.
  \item \textbf{El efecto retenido no es conjuntamente seguro.} $+0.066$ de ReT medio, pero con costo
        de peor-producto y no-resueltos; no es desplegable tal cual.
  \item \textbf{No hay óptimo de control certificado.} El MPC es un controlador aproximado de horizonte
        deslizante; que su creencia sea óptima para \emph{inferencia} no prueba que sea óptimo para
        \emph{control}.
  \item \textbf{Validación física pendiente.} Cualquier extensión de riesgo acoplado a producto
        requiere validación escrita de Garrido antes de ejecutarse.
\end{itemize}

%==============================================================
\section{Resumen ejecutivo en una página}
\begin{itemize}[nosep]
  \item Reconstruimos, con fidelidad exacta (0/47\,546), una cadena de suministro militar de 13
        operaciones y su métrica de resiliencia ReT.
  \item Construimos un marco de falsación de 4 niveles con una \textbf{frontera de lazo abierto exacta
        de 65\,536 planes} como clave de respuestas.
  \item \textbf{Logro 1 (Submission A):} el RL aprende control de lazo cerrado real —le gana a las
        65\,536 políticas fijas ($+0.07$ a $+0.12$)— pero \textbf{empata al MPC} ($\Delta_N{\approx}0$,
        equivalencia TOST, robusta a 5 pruebas); y el MPC es además ${\sim}7\times$ más rápido. Entrega
        el puente lazo-abierto$\to$lazo-cerrado de Garrido y aporta un protocolo de decisión ``cuándo
        entrenar''.
  \item \textbf{Logro 2 (Q-R1):} la retención de conocimiento entre campañas \textbf{mejora causalmente
        la resiliencia de arranque en frío} ($+0.066$, LCB $+0.052$, dosis-respuesta perfecta,
        anti-oráculo correcto), confirmada prospectivamente en seeds selladas —el aprendizaje acumulado
        de Garrido. \emph{Pero} aún no es conjuntamente segura (costo de equidad de producto).
  \item \textbf{Estrella polar restante:} un híbrido learning-augmented MPC que capture esa ganancia
        \emph{sin} el costo de servicio, o bien un certificado honesto de frontera. La disciplina —seeds
        selladas, verificación adversarial, métrica intacta— es lo que hace creíbles tanto los ``sí''
        como los ``no''.
\end{itemize}

\vfill
\noindent\rule{\linewidth}{0.4pt}\\
{\footnotesize Documento acompañante generado el 22 de julio de 2026. Los números provienen de
artefactos con custodia SHA-256: la tabla de Submission A (\code{source\_of\_truth.json}, veredicto
\code{STOP\_Q\_NO\_REPLICATED\_LEARNED\_ADAPTATION}), el benchmark de latencia
(\code{latency\_benchmark\_v1}) y la adjudicación verificada del confirmatorio Q-R1
(\code{STOP\_REPAIRED\_Q\_R1\_NO\_RETAINED\_INFORMATION\_PASS}, verificación A11 independiente sobre los
outputs crudos del VPS). La métrica ReT es fiel a la tesis Garrido-Ríos 2017.}

\end{document}
